\documentclass[11pt]{article}
\usepackage[margin=1in]{geometry}
\usepackage{amsmath,amssymb}
\usepackage{booktabs}
\usepackage{hyperref}
\usepackage{xcolor}
\usepackage{tcolorbox}

\title{Independent Replication Report:\\
Orbital Hall Effect in Spin-3/2 Hole-Doped Semiconductors (Ge)\\
\large Cullen, Wang \& Culcer, arXiv:2509.20436v3 (2025/2026)}
\author{REPLICATE-PROJECT / TEXTURES-100 --- conventional-Kubo rebuild}
\date{2026-07-19}

\begin{document}
\maketitle

\begin{tcolorbox}[colback=yellow!8!white,colframe=black!60]
\textbf{VERDICT: PARTIAL} \quad (Coverage 5/10, Agreement 5/10)\\[2pt]
\textbf{Headline computed (this work):} $\sigma^{L_z}_{xy,\mathrm{conv}} \approx 49$ $(\hbar/e)\,\Omega^{-1}\mathrm{cm}^{-1}$
(conventional interband Kubo term only; converged, $<3\%$ drift over $N=21/31/41$).\\
\textbf{Paper headline:} $\sigma_{\mathrm{OHE}} \sim 10^{3}$ $(\hbar/e)\,\Omega^{-1}\mathrm{cm}^{-1}$
(TOTAL: conventional $+$ quantum corrections).\\[2pt]
The $\sim\!20\times$ gap is \emph{not} a disagreement. The paper's Fig.~2 states explicitly
that the quantum corrections $\Delta\sigma_1,\Delta\sigma_2$ \emph{dominate} the total and are what
push it to $\sim\!10^3$. Our conventional-only rebuild is the correct value of the one term we built;
the headline requires the unbuilt quantum-correction terms of the paper's Eq.~(6).
\end{tcolorbox}

\section{The paper and its central claim}
Cullen \emph{et al.} propose bulk $p$-type semiconductors (focus: Ge; also Si, GaAs, InAs, InSb)
as orbitronic platforms. Using the modern theory of orbital magnetisation applied to the
Luttinger--Kohn--Bir--Pikus valence-band Hamiltonian, and \emph{incorporating recently-discovered
quantum corrections to the orbital current}, they report an orbital Hall conductivity of order
$10^{3}\,(\hbar/e)\,\Omega^{-1}\mathrm{cm}^{-1}$ --- exceeding the spin Hall effect by 2--3 orders
of magnitude. The one comparison anchor for this replication is that $\sim\!10^3$ headline for Ge,
and its decomposition into a conventional piece plus dominant quantum corrections (paper Fig.~2).

\paragraph{Numbered claims checked.}
\begin{enumerate}
\item[C1.] The 4$\times$4 spherical Luttinger model for Ge has \emph{finite} Berry curvature and OAM
even under inversion symmetry, of the monopole-like form $\Omega^m_{i,\mathbf k}=-(k_i/k^3)(J_z)_{mm}$ (Eq.~3).
\item[C2.] The orbital Hall conductivity is of the right order for orbitronics; the \emph{total} is
$\sim\!10^3\,(\hbar/e)\,\Omega^{-1}\mathrm{cm}^{-1}$ (abstract, Fig.~2).
\item[C3.] The quantum corrections $\Delta\sigma_1,\Delta\sigma_2$ are the \emph{dominant} contributions;
the conventional interband piece $\sigma_{\mathrm{conv}}$ is sub-dominant (Fig.~2, Sec.~text).
\end{enumerate}

\section{Governing equations reimplemented}
The full model is the 6$\times$6 LKBP Hamiltonian, Eq.~(1). For Ge (large split-off gap $\delta_{so}$)
the upper 4$\times$4 block suffices, and we use the \textbf{spherical approximation}, Eq.~(2):
\begin{equation}
H_0 = -\frac{\hbar^{2}}{2m_0}\Big[(\gamma_1+\tfrac52\bar\gamma)\,k^2 - 2\bar\gamma\,(\mathbf{k}\cdot\mathbf{J})^2\Big],
\qquad \bar\gamma=\tfrac12(\gamma_2+\gamma_3),
\end{equation}
with $\mathbf{J}$ the spin-3/2 matrices (heavy holes $J_z=\pm3/2$, light holes $J_z=\pm1/2$).
For Ge we use $\gamma_1=13.38$, $\gamma_2=4.25$, $\gamma_3=5.69 \Rightarrow \bar\gamma=4.97$.

The Berry curvature (Eq.~3) and equilibrium OAM (Eq.~4) are analytic:
\begin{equation}
\Omega^{m}_{i,\mathbf k}=-\frac{k_i}{k^3}(J_z)_{mm},\qquad
L^{m}_{i,\mathbf k}=3\hbar\bar\gamma\,\frac{k_i}{\cdots}(J_z)_{mm}.
\end{equation}

The orbital current with all matrix elements is Eq.~(6),
$\langle J^\alpha_\delta\rangle=\mathrm{Tr}\big[\tfrac12\{\hat L^\alpha,\hat v_\delta\}\,\rho\big]$.
The \emph{conventional} piece keeps only off-diagonal velocity and band-diagonal OAM matrix
elements. We implement exactly this conventional interband (intrinsic Kubo / Berry-curvature)
term, with the orbital current operator $j^{L_z}_x=\tfrac12\{L_z,v_x\}$:
\begin{equation}
\sigma^{L_z}_{xy,\mathrm{conv}}
= \frac{e^2}{\hbar}\,\frac{1}{(2\pi)^3}\int d^3k \sum_{n\,\mathrm{occ}}\sum_{m\neq n}
\frac{2\hbar^2\,\mathrm{Im}\big[(j^{L_z}_x)_{nm}(v_y)_{mn}\big]}{(\varepsilon_n-\varepsilon_m)^2}.
\end{equation}
The velocity operator is $v_i=\tfrac{1}{\hbar}\partial H_0/\partial k_i$, obtained analytically from
Eq.~(2). We do \textbf{not} build the quantum corrections $\Delta j_{1,2,3}$ of Eq.~(6) --- this is
the deliberate scope of the replication and the source of the coverage cap.

\section{Method}
Direct numerical Kubo evaluation on a uniform 3D $k$-grid, spherical 4$\times$4 model:
\begin{itemize}
\item Build $H_0(\mathbf k)$ and its $\mathbf k$-gradients analytically; \texttt{eigh} at each grid point.
\item Rotate $v_x,v_y$ and $L_z$ into the eigenbasis; form $j^{L_z}_x=\tfrac12\{L_z,v_x\}$.
\item Sum the interband Berry-curvature-weighted term over occupied hole states
($\varepsilon_n\le E_F$), degenerate pairs ($|\varepsilon_n-\varepsilon_m|<10^{-24}$ J) masked.
\item Normalise $\int d^3k/(2\pi)^3$; convert SI $\to$ S/cm $\equiv (\hbar/e)\,\Omega^{-1}\mathrm{cm}^{-1}$.
\end{itemize}
Grid: cubic box $k_{\max}=4\times10^8\,\mathrm{m}^{-1}$, $N^3$ points ($N=21,31,41$),
$E_F=10$ meV (also 5 meV). Fermi-surface coverage checked: the heavy-hole Fermi wavevector
$k_F^{hh}=\sqrt{2m_0 E_F/[\hbar^2(\gamma_1-2\bar\gamma)]}=2.76\times10^8\,\mathrm{m}^{-1}<k_{\max}$,
so the grid encloses the full occupied hole pocket.

\section{Results}
\begin{table}[h]\centering
\begin{tabular}{lcccc}
\toprule
Run & $N$ & $E_F$ (meV) & $\sigma^{L_z}_{xy,\mathrm{conv}}$ $(\hbar/e)\Omega^{-1}\mathrm{cm}^{-1}$ & $t$ (s)\\
\midrule
coarse & 21 & 10 & 48.30 & 0.2\\
med    & 31 & 10 & 49.92 & 0.9\\
fine   & 41 & 10 & 49.41 & 2.4\\
fine   & 41 & 5  & 33.77 & 3.8\\
\bottomrule
\end{tabular}
\caption{Grid-convergence of the conventional interband OHE for Ge. The $N=21/31/41$ values
(48.3 / 49.9 / 49.4) drift $<3\%$ $\Rightarrow$ \textbf{converged}. Live re-run (2026-07-19,
\texttt{/home/stevens/comfyui-env/bin/python}) reproduced $N=31$ to the quoted digits (49.9217).}
\end{table}

\begin{table}[h]\centering
\begin{tabular}{lccc}
\toprule
Quantity & This work & Paper & Status\\
\midrule
$\sigma^{L_z}_{xy,\mathrm{conv}}$ (Ge) & $\approx 49$ & sub-dominant term in Fig.~2 & consistent (right order for $\sigma_{\mathrm{conv}}$)\\
$\sigma_{\mathrm{OHE}}$ total (Ge)     & not built    & $\sim 10^3$ & gap = unbuilt $\Delta\sigma_{1,2}$\\
Finite $\Omega,L$ under inversion (C1) & yes, $\Omega_i\propto -k_i/k^3\,J_z$ & yes (Eq.~3) & \textbf{reproduced}\\
$\Delta\sigma_{1,2}$ dominate (C3)     & not built    & yes (Fig.~2) & consistent (our $\sigma_{\mathrm{conv}}\ll10^3$)\\
\bottomrule
\end{tabular}
\caption{This-work vs paper. The conventional term reproduces the paper's physics and order for that
term; the total headline requires the quantum corrections, which are out of scope.}
\end{table}

\section{Comparison and verdict}
Our converged conventional-only conductivity, $\sigma^{L_z}_{xy,\mathrm{conv}}\approx49\,
(\hbar/e)\Omega^{-1}\mathrm{cm}^{-1}$, sits $\sim\!20\times$ below the paper's $\sim\!10^3$ headline.
This is \emph{consistent}, not contradictory: the paper's central methodological point (abstract;
Fig.~2; Sec. text lines 590--591 ``the quantum corrections $\Delta\sigma_{1,2}$ are the dominant
contributions'') is precisely that the conventional interband term is sub-dominant and the total is
carried by the quantum corrections $\Delta j_{1},\Delta j_{2}$ of their Eq.~(6). A tight
conventional-only rebuild landing 1--2 orders below the total is the expected signature of that claim.
The core single-particle physics --- spherical Luttinger bands, finite Berry curvature/OAM under
inversion symmetry of the form $\Omega_i=-(k_i/k^3)J_z$, and the right order of magnitude for the
conventional piece --- \textbf{reproduces}. \textbf{Verdict: PARTIAL.}

\section{Gaps, conventions and scope}
\begin{enumerate}
\item \textbf{Quantum corrections not built.} $\Delta\sigma_1$ (interband-polarization / dipole
rotation) and $\Delta\sigma_2$ (interband OAM matrix elements) --- the paper's \emph{dominant} terms
--- are not implemented; nor is $\Delta\sigma_3$ (the $[\hat r,\hat v]$ term, opposite sign). Closing
the headline requires the full Eq.~(6) covariant-derivative machinery of the modern theory of orbital
magnetisation. This is the coverage cap, and it is a deliberate scope choice, not a physics failure.
\item \textbf{Spherical 4$\times$4 approximation only.} Full 4$\times$4 (cubic anisotropy) and 6$\times$6
(split-off band) models, used by the paper for numerical accuracy, are not built. Adequate for Ge
per the paper (large $\delta_{so}$), but caps quantitative fidelity.
\item \textbf{Conventional-Kubo convention.} We adopt the community-standard Go-\emph{et al.} form
(charge-Hall Kubo prefactor $e^2/\hbar$, one velocity replaced by $j^{L_z}_x$, $L_z$ dimensionless).
The paper's ``proper'' (Eq.~7-analogue) current with the tilde/check degenerate/non-degenerate
distinction is not separately reconstructed.
\item \textbf{Grid \& Fermi surface (verified, not a gap).} $k_{\max}>k_F^{hh}$ confirmed; convergence
$<3\%$ over $N=21/31/41$. The number is trustworthy \emph{for the term computed}.
\item \textbf{Extraction tooling degraded (not physics).} \texttt{marker}/\texttt{nougat} absent;
extraction/*.\{md,mmd\} are honest \texttt{pdftotext} interims with hand-transcribed equations.
\end{enumerate}

\section{Reproducibility}
Interpreter: \texttt{/home/stevens/comfyui-env/bin/python} (numpy 2.3.5). Driver:
\texttt{report/evidence/ohe\_spherical.py}. Evidence JSON:
\texttt{report/evidence/cullen2025\_result.json}.
\begin{verbatim}
cd /home/stevens/textures-100/corpus/textures-orbital-cullen2025/
/home/stevens/comfyui-env/bin/python work/ohe_spherical.py
# -> writes work/cullen2025_result.json with the N21/31/41 convergence
#    table + _verdict block (grid_covers_FS, computed value, PARTIAL).
\end{verbatim}
This \texttt{.tex} ships as source; no LaTeX engine is required on the packaging host and it compiles
off-host with \texttt{pdflatex} (packages: amsmath, booktabs, hyperref, tcolorbox).

\end{document}
